\documentclass[letterpaper,11pt]{article}

\usepackage[letterpaper, left=0.8in, right=0.8in, top=0.7in, bottom=0.7in]{geometry}
\usepackage[hidelinks]{hyperref}
\usepackage{parskip}
\usepackage{setspace}

\setstretch{1.04}
\pagestyle{empty}

\begin{document}

\begin{center}
    {\Large \textbf{Aryan Miriyala}}\\
    \vspace{2pt}
    +1 419-315-0444 \textbar{} \href{mailto:aryanmiriyala@gmail.com}{aryanmiriyala@gmail.com} \textbar{} \href{https://www.linkedin.com/in/aryan-miriyala-7788a21b6/}{linkedin.com/in/aryan-miriyala} \textbar{} \href{https://github.com/aryanmiriyala}{github.com/aryanmiriyala}
\end{center}

\vspace{10pt}

Dear Runpod Hiring Team,

I am applying for the Software Engineer (Full-Stack) role because the product sits in a technical area I want to build for directly: practical tools that make AI infrastructure usable for the developers and researchers trying to ship real workloads. Runpod's focus on APIs, SDKs, containerized inference, Serverless GPU endpoints, and developer experience is a strong match for the way I have been building full-stack AI and cloud-adjacent systems across internships, research, and projects.

At SmartSolve, I have been working on secure internal software with Next.js, TypeScript, PostgreSQL, Drizzle ORM, SSO, auth middleware, Docker, Git, and AI-assisted development workflows. One project is a full-stack onboarding tracker for HR workflows; another is an AI-enabled QMS dependency register where the challenge was turning ambiguous document relationships into an implementable system plan. That work maps well to Runpod's need for engineers who can move between user interfaces, APIs, data models, debugging, and collaboration with technical and non-technical stakeholders.

My project work is especially relevant to tools for AI engineers. I built FalconGraph Search as a source-grounded campus answer engine with a C++/OpenMP crawler, Python document processing, embeddings, FAISS retrieval, a FastAPI RAG endpoint, and a Next.js interface that returns cited answers with source context. I also built Fix-It-Flow, a Next.js 14/TypeScript PWA with typed API routes, camera input, voice commands, Gemini Vision, LLM reasoning, ElevenLabs text-to-speech, and DynamoDB-backed session persistence. In both cases, the useful part was not just calling a model; it was designing the surrounding API boundaries, state, retrieval, UI feedback, and error-prone workflow details so the system stayed understandable.

I would bring Runpod a practical early-career full-stack profile: Python and TypeScript, React/Next.js, Express/FastAPI-style API work, Docker, Git, AWS data and event workflows, and a strong interest in developer-facing infrastructure. I am excited by the chance to learn from senior engineers while contributing to tools that make AI infrastructure easier to integrate, test, and use.

\vspace{8pt}

Best regards,\\
Aryan Miriyala

\end{document}

